
\subsection{Affine flag varieties and jet geometry}

\begin{definition}[Affine flag variety]\label{def:aff-flag}
The affine flag variety is:
\[
\mathrm{Fl}_{\mathrm{aff}} = G(\mathbb{C}((t)))/I
\]
where $G(\mathbb{C}((t)))$ is the loop group and
$I = \{g \in G(\mathbb{C}[[t]]) : g \bmod t \in B(\mathbb{C})\}$ is the Iwahori subgroup
(the preimage of $B$ under evaluation at $t = 0$; strictly larger than $B(\mathbb{C}[[t]])$).
\end{definition}

\begin{theorem}[Jet bundle realization \cite{BD04}; \ClaimStatusProvedElsewhere]\label{thm:jet-flag}
There is a natural isomorphism:
\[
J^{\infty}(G/B) \cong \mathrm{Fl}_{\mathrm{aff}}^{\mathrm{thick}}
\]
where the right side is the ``thick'' flag variety with formal neighborhood structure.

This identifies:
\begin{itemize}
\item Points in $G/B$: Finite-dimensional flags
\item Jets at a point: Formal deformations of flags
\item W-algebra generators: Functions on jet space
\end{itemize}
\end{theorem}

\begin{proof}[Construction via loop spaces]
\emph{Step 1: Loop Space Interpretation.} The affine Grassmannian parametrizes lattices in the loop space:
\[
\mathrm{Gr}_G = G(\mathbb{C}((t)))/G(\mathbb{C}[[t]]) = \{\text{lattices } L \subset \mathbb{C}((t))^n : t^N \cdot L_{\mathrm{std}} \subset L \subset t^{-N} \cdot L_{\mathrm{std}} \text{ for some } N\}
\]
The affine flag variety $\mathrm{Fl}_{\mathrm{aff}} = G(\mathbb{C}((t)))/I$ (where $I$ is an Iwahori subgroup) parametrizes complete \emph{chains} of lattices, not individual lattices.

\emph{Step 2: Jet Interpretation.} A jet of a flag at a point corresponds to:
\begin{itemize}
\item Order 0: The flag itself
\item Order 1: Infinitesimal deformation
\item Order $k$: $k$-th order Taylor expansion
\end{itemize}

\emph{Step 3: Dictionary.} Under the isomorphism:
\begin{align}
\text{W-algebra generator } W^{(s)} &\leftrightarrow \text{Function on jets of order } s-1 \\
\text{OPE singularity} &\leftrightarrow \text{Poisson bracket on jet bundle} \\
\text{Normal ordering} &\leftrightarrow \text{Weyl quantization}
\end{align}
\end{proof}

\subsection{Chiral de Rham complex and resolution}

\begin{theorem}[\texorpdfstring{$\mathcal{W}$}{W}-algebras as chiral de Rham \cite{Arakawa}; \ClaimStatusProvedElsewhere]\label{thm:w-cdr}
At critical level $k = -h^\vee$, the W-algebra admits a resolution:
\[
\mathcal{W}^{-h^\vee}(\mathfrak{g}, e) \cong H^*_{\mathrm{DS}}(\Omega^{\mathrm{ch}}_{G/P_e})
\]
where:
\begin{itemize}
\item $\Omega^{\mathrm{ch}}_{G/P_e}$ is the chiral de Rham complex of the partial flag variety
\item $H^*_{\mathrm{DS}}$ is Drinfeld--Sokolov cohomology
\item $P_e$ is the parabolic determined by $e$
\end{itemize}
\end{theorem}

\begin{proof}
\emph{Step 1: Factorization algebra model.} The chiral de Rham complex defines a factorization algebra $U \mapsto \Omega^{\mathrm{ch}}(U \times_{X} G/P_e)$ for open $U \subset X$.

\emph{Step 2: Drinfeld--Sokolov reduction.} The DS reduction is implemented by the BRST differential $Q_{\mathrm{DS}} = \{Q_{\mathrm{BRST}}, -\}$, with screening charges given by integrals of nilpotent currents associated to simple roots. The W-algebra generators emerge as $Q_{\mathrm{DS}}$-closed elements not in the image of $Q_{\mathrm{DS}}$.

\emph{Step 3: Bar complex identification.} By the factorization algebra
formulation, the bar complex computes
\[
\bar{B}^{\mathrm{ch}}(\mathcal{W}^{-h^\vee}(\mathfrak{g}, e))
\cong \mathrm{C}_*(\mathrm{Maps}(X, G/P_e)).
\]
This identifies the bar differential with the geometric boundary operator on
the mapping-space chain model for holomorphic maps into $G/P_e$.
\end{proof}

\subsection{Explicit bar complex coalgebra for $\mathcal{W}$-algebras}

\begin{theorem}[\texorpdfstring{$\mathcal{W}$}{W}-algebra bar coalgebra; \ClaimStatusProvedHere]\label{thm:w-bar-coalg}
For $\mathcal{W}^k(\mathfrak{g}, e)$, the bar complex carries:
\begin{enumerate}
\item \emph{Comultiplication.} For generators $W^{(s_i)}$ at points $z_i$:
\[
\Delta(W^{(s_1)}_{z_1} \otimes \cdots \otimes W^{(s_n)}_{z_n}) = 
\sum_{\text{partitions}} \mathrm{Fusion}_I \otimes \mathrm{Fusion}_J
\]
where fusion uses the W-algebra OPE algebra.

\item The coalgebra structure is dual to the intersection pairing:
\[
\langle \alpha, \beta \rangle = \int_{G/P_e} \alpha \wedge \beta
\]
the intersection pairing on the flag variety cohomology.

\item At non-critical level, the comultiplication acquires quantum corrections:
\[
\Delta_{\hbar}(W) = \Delta_0(W) + \hbar \Delta_1(W) + \hbar^2 \Delta_2(W) + \cdots
\]
where $\hbar = 1/(k + h^\vee)$ is the quantum parameter.
\end{enumerate}
\end{theorem}

\begin{proof}
\emph{(1) Comultiplication.}  The chiral bar complex $\bar{B}^{\mathrm{ch}}(\mathcal{W}^k)$ is a coalgebra with factorization coproduct.  For generators $W^{(s_i)}$ placed at points $z_i \in C_n(X)$, the coproduct decomposes the configuration into two sub-configurations:
\[
\Delta\bigl(W^{(s_1)}_{z_1} \otimes \cdots \otimes W^{(s_n)}_{z_n}\bigr) = \sum_{I \sqcup J = \{1,\ldots,n\}} \pm\, \bigl(W^{(s_i)}_{z_i}\bigr)_{i \in I} \otimes \bigl(W^{(s_j)}_{z_j}\bigr)_{j \in J}
\]
where the sum is over all partitions of the point set into two disjoint subsets (reflecting the factorization structure on configuration spaces, \emph{not} the deconcatenation coproduct of the algebraic bar complex, which would sum only over consecutive splittings).  Signs are determined by Koszul conventions.

\emph{(2) Intersection pairing.}  At critical level $k = -h^\vee$, the bar complex computes chains on maps into $G/P_e$ (Theorem~\ref{thm:w-cdr}).  The coalgebra structure is Poincar\'e dual to the intersection product on $H^*(G/P_e)$: for $\alpha, \beta \in H^*(G/P_e)$, $\langle \alpha, \beta \rangle = \int_{G/P_e} \alpha \wedge \beta$.

\emph{(3) Quantum corrections.}  At general level, the coproduct acquires corrections from the curvature.  Setting $\hbar = 1/(k + h^\vee)$, the corrected coproduct $\Delta_\hbar = \Delta_0 + \hbar \Delta_1 + \cdots$ satisfies the co-$A_\infty$ relations dual to those of Theorem~\ref{thm:w-ainfty-ops}.  Each $\Delta_n$ is determined by integrating over the codimension-$n$ strata of $\overline{C}_*(X)$.
\end{proof}

\begin{example}[Virasoro bar complex coproduct]
The bar complex $\bar{B}(\mathrm{Vir}_c)$ carries the partition coproduct (sum over all set-partitions $I \sqcup J$, not merely consecutive splittings). On generators $[L_n] \in \bar{B}^1$:
\begin{align}
\Delta([L_n]) &= [L_n] \otimes 1 + 1 \otimes [L_n] \quad \text{(primitive)}\\
\Delta(c) &= c \otimes 1 + 1 \otimes c \quad \text{(primitive)}
\end{align}
The non-trivial coalgebra structure appears in higher bar degrees, where $\Delta([L_m | L_{n-m}]) = [L_m] \otimes [L_{n-m}] + \cdots$ encodes the OPE data.
\end{example}

\subsection{$\mathcal{W}$-algebras at general levels and screening operators}

\begin{definition}[Screening operators]\label{def:screening}
For $\mathcal{W}^k(\mathfrak{g}, e)$, screening operators are:
\[
S_{\alpha} = \oint_{\gamma_\alpha} V_\alpha(z) \, dz
\]
where:
\begin{itemize}
\item $V_\alpha$ are vertex operators of specific weights
\item $\gamma_\alpha$ are cycles in $H_1(X \setminus \{\text{poles}\})$
\item $[Q_{\mathrm{BRST}}, S_\alpha] = 0$ (screening condition)
\end{itemize}
\end{definition}

\begin{theorem}[Screening resolution \cite{FF-wakimoto}; \ClaimStatusProvedElsewhere]\label{thm:screen-res}
The W-algebra admits a free field resolution:
\[
\mathcal{W}^k(\mathfrak{g}, e) = \mathrm{Ker}(S_1, \ldots, S_r : \mathcal{FF} \to \mathcal{FF}')
\]
where $\mathcal{FF}$ is a free field algebra (Wakimoto module).
\end{theorem}

\subsection{Non-principal nilpotents and the DS hierarchy}
\label{sec:non-principal-ds}
\index{Drinfeld--Sokolov reduction!non-principal nilpotent}
\index{nilpotent orbit!W-algebra}

The DS reduction $\mathcal{W}^k(\mathfrak{g},f)$ depends on the nilpotent orbit of $f$.
The zero orbit gives the universal affine algebra, while the principal orbit gives the principal W-algebra.
Other orbits are ordered between these in the nilpotent orbit closure relation.
This orbit order does not provide vertex embeddings between their reductions.

\begin{remark}[Reduction and vertex inclusions]
The nilpotent orbit determines a reduction complex and its degree-zero cohomology.
It does not determine an inclusion into the bare affine algebra.
At level zero, such an inclusion for subregular $\mathfrak{sl}_3$ contradicts the vacuum product $J_{(1)}J=\mathbf1$
(Proposition~\ref{prop:cbc28-nonembedding}).
The strict reduction map instead lands in the affine--ghost complex.
\end{remark}

\begin{computation}[DS hierarchy for \texorpdfstring{$\mathfrak{sl}_3$}{sl_3}]
\label{comp:sl3-ds-hierarchy}
\index{sl3@$\mathfrak{sl}_3$!nilpotent orbits}
The nilpotent orbits of $\mathfrak{sl}_3$ are:
\begin{center}
\renewcommand{\arraystretch}{1.2}
\begin{tabular}{lcccl}
\textbf{Orbit} & \textbf{Partition} & \textbf{$\dim \mathbb{O}$} &
\textbf{W-algebra} & \textbf{Generators} \\
\hline
$\{0\}$ & $(1,1,1)$ & $0$ & $\widehat{\mathfrak{sl}}_{3,k}$ &
$J^a$ ($a = 1,\ldots,8$) \\
$\mathbb{O}_{\min}$ & $(2,1)$ & $4$ &
$\mathcal{W}^k(\mathfrak{sl}_3, f_{\min})$ &
$J,T,G^+,G^-$ \\
$\mathbb{O}_{\mathrm{reg}}$ & $(3)$ & $6$ &
$\mathcal{W}^k(\mathfrak{sl}_3, f_{\mathrm{prin}}) = \mathcal{W}_3$ &
$T$ (spin 2) + $W$ (spin 3)
\end{tabular}
\end{center}

The minimal nonzero orbit is also subregular.
Its reduction has four even strong generators $J,T,G^+,G^-$, of weights $1,2,3/2,3/2$ for $k\ne-3$.
Their explicit representatives are~\eqref{eq:cbc27-generators} and~\eqref{eq:cbc27-bp-fields}.
The centralizer in grading degree zero $\mathfrak{sl}_3^f\cap\mathfrak g_0=\mathbb Cj\oplus\mathbb Cu$
uses the integral good grading and has brackets $[j,u]=u$.
For the symmetric conformal vector $T$, the weight-one subspace is the single Heisenberg current $J$.
Its coefficient is $(2k+3)/3$.

Theorem~\ref{thm:cbc27-subregular-map} constructs a strict unital map into the full affine--ghost complex and proves it is a quasi-isomorphism.
Proposition~\ref{prop:cbc27-binary} constructs its meromorphic binary collision map.
In the PBW associated graded, the differential polynomial generators have weights $1,2,3/2,3/2$.
Consequently the conformal vacuum series is
\[
 \prod_{r\geq0}\frac{1}{(1-q^{r+1})(1-q^{r+2})(1-q^{r+3/2})^2}.
\]
The three lowest nonconstant weight spaces have dimensions $1,2,3$ at weights $1,3/2,2$.
These are vertex-state dimensions.
A tensor length in a chiral bar also contains derivatives and configuration-space coefficients, so these numbers do not count its chain groups.
The reflected orbit $(2,1)$ is unchanged, while the conformal charge calculation is Corollary~\ref{cor:cbc27-bp-levels}.
\end{computation}

\begin{computation}[Miura transform for \texorpdfstring{$\mathcal{W}_3$}{W_3}]
\label{comp:miura-w3}
\index{Miura transform!$\mathcal{W}_3$}

The Miura transform realizes $\mathcal{W}_3$ generators inside the
Heisenberg algebra $\mathcal{H}^{\otimes 2}$ (two free bosons $\phi_1, \phi_2$
for the rank-2 Cartan of $\mathfrak{sl}_3$).

\emph{Differential operator.}
The Miura operator is the $3 \times 3$ matrix differential operator:
\begin{equation}\label{eq:miura-sl3}
L = (\partial + \alpha_1(\partial\phi))
(\partial + \alpha_2(\partial\phi))
(\partial + \alpha_3(\partial\phi))
\end{equation}
where $\alpha_1 = (1, 0)$, $\alpha_2 = (-1/2, \sqrt{3}/2)$,
$\alpha_3 = (-1/2, -\sqrt{3}/2)$ are the weights of the fundamental
representation of $\mathfrak{sl}_3$ (satisfying $\alpha_1+\alpha_2+\alpha_3=0$),
and $\partial\phi = (\partial\phi_1, \partial\phi_2)$.

\emph{Expansion.}
Setting $a_i = \alpha_i(\partial\phi) = \alpha_i^1 \partial\phi_1
+ \alpha_i^2 \partial\phi_2$:
\begin{align*}
L &= \partial^3 + (a_1+a_2+a_3)\partial^2
+ (a_1 a_2 + a_1 a_3 + a_2 a_3 + a_1' + a_2')\partial \\
&\quad + (a_1 a_2 a_3 + a_1 a_2' + a_1' a_3)
\end{align*}
Since $a_1 + a_2 + a_3 = (\alpha_1+\alpha_2+\alpha_3)(\partial\phi) = 0$,
the $\partial^2$ coefficient vanishes identically.

\emph{W-algebra generators.}
The coefficients of $L$ (after removing the $\partial^2$ term) are:
\begin{align}
T &= -(a_1 a_2 + a_1 a_3 + a_2 a_3 + a_1' + a_2')
\label{eq:miura-T} \\
W &= a_1 a_2 a_3 + a_1 a_2' + a_1' a_3
\label{eq:miura-W}
\end{align}

\emph{Explicit formulas in terms of $\partial\phi_1, \partial\phi_2$.}
Using $a_1 = \partial\phi_1$,
$a_2 = -\frac{1}{2}\partial\phi_1 + \frac{\sqrt{3}}{2}\partial\phi_2$,
$a_3 = -\frac{1}{2}\partial\phi_1 - \frac{\sqrt{3}}{2}\partial\phi_2$:
\begin{align}
T &= -\frac{1}{2}\bigl((\partial\phi_1)^2 + (\partial\phi_2)^2\bigr)
- \frac{1}{\sqrt{3}}\partial^2\phi_2
\label{eq:miura-T-explicit} \\
W &= \frac{1}{3\sqrt{3}}\bigl((\partial\phi_2)^3
- 3(\partial\phi_1)^2 \partial\phi_2\bigr)
+ \text{(derivative terms)}
\label{eq:miura-W-explicit}
\end{align}
The stress tensor $T$ is the Sugawara tensor for two free bosons
(central charge $c = 2$ at $k \to \infty$) plus an improvement term
$\partial^2\phi_2$ (the ``background charge'') controlled by the level.

At level $k$, the full Miura transform includes the background charge
$Q = (k+3)^{-1/2}$:
\begin{equation}
T_k = -\frac{1}{2}(\partial\phi)^2 + Q \rho \cdot \partial^2\phi
\end{equation}
where $\rho = (\alpha_1 + \alpha_2)/2$ is the Weyl vector.  The
central charge becomes $c_k = 2 - 24Q^2|\rho|^2 = 2 - 24(k+2)^2/(k+3)$.

The background charge $Q$ is the Miura manifestation of the
universal MC class: $\kappa(\mathcal{W}_3^k) = c_k/2$, and the
Koszul duality $k \mapsto -k-6$ sends $Q \mapsto -Q$, implementing
$\Theta_{\mathcal{W}_3^k} + \Theta_{\mathcal{W}_3^{-k-6}} = 0$
(Theorem~\ref{thm:explicit-theta}).  The Miura embedding
$\mathcal{W}_3^k \hookrightarrow \mathcal{H}^{\otimes 2}_Q$
is the free-field avatar of the fact that $\Theta$ factors
through the Heisenberg curvature $\kappa(\mathcal{H}_Q) = Q^2/2$.

\emph{Koszul duality via Miura.}
The Miura transform exhibits the free field embedding
$\mathcal{W}_3^k \hookrightarrow \mathcal{H}^{\otimes 2}_{Q}$.
Under Koszul duality, the background charge transforms as
$Q \mapsto -Q$ (since $k \mapsto -k-6$ gives
$Q' = (-k-6+3)^{-1/2} = (-(k+3))^{-1/2}$), and the screening
charges $S_i = \oint e^{\alpha_i \cdot \phi/Q}$ dualize to
$S_i' = \oint e^{-\alpha_i \cdot \phi/Q}$.
\end{computation}


\subsection{Connection to integrable systems}

\begin{theorem}[\texorpdfstring{$\mathcal{W}$}{W}-algebras and Toda systems \cite{FF}; \ClaimStatusProvedElsewhere]\label{thm:w-toda}
The W-algebra $\mathcal{W}^k(\mathfrak{g}, e_{\mathrm{prin}})$ at the principal nilpotent is equivalent to:
\[
\mathcal{W}^k(\mathfrak{g}, e_{\mathrm{prin}}) \cong \text{Quantization of Toda system for } \mathfrak{g}
\]
where the Toda system is the integrable system with Lax operator:
\[
L = \partial + f_{\mathrm{prin}} + \sum_{i} \phi_i h_i
\]
\end{theorem}

\begin{computation}[Toda Hamiltonians for \texorpdfstring{$\mathfrak{sl}_3$}{sl_3}]
\label{comp:toda-sl3}
\index{Toda system!sl3@$\mathfrak{sl}_3$}

The $\mathfrak{sl}_3$ Toda system has two fields $\phi_1, \phi_2$
(corresponding to the rank-2 Cartan) and Lax operator:
\begin{equation}
L = \begin{pmatrix}
\partial + \partial\phi_1 & 0 & 0 \\
e^{\phi_2 - \phi_1} & \partial + \partial\phi_2 - \partial\phi_1 & 0 \\
0 & e^{-\phi_2} & \partial - \partial\phi_2
\end{pmatrix}
\end{equation}
(the lower-triangular form reflects the principal nilpotent
$f_{\mathrm{prin}} = F_1 + F_2$).

\emph{Classical Hamiltonians.}
The conserved charges are $H_s = \operatorname{Res} \operatorname{tr}(L^s)$
for $s = 2, 3$:
\begin{align}
H_2 &= \int \Bigl(\frac{1}{2}(\partial\phi_1)^2
+ \frac{1}{2}(\partial\phi_2)^2
- \partial\phi_1 \partial\phi_2
+ e^{\phi_2 - \phi_1} + e^{-\phi_2}\Bigr)\, dz \\
H_3 &= \int \Bigl((\partial\phi_1)(\partial\phi_2)(\partial\phi_2 - \partial\phi_1)
+ \text{exponential terms}\Bigr)\, dz
\end{align}
$H_2$ is the Toda Hamiltonian (the energy); $H_3$ is the cubic
integral of motion.  In Poisson involution: $\{H_2, H_3\} = 0$.

\emph{Quantization.}
The quantum Toda Hamiltonians are the $\mathcal{W}_3$ generators:
$T \leftrightarrow H_2$ and $W \leftrightarrow H_3$ under the Miura
transform (Computation~\ref{comp:miura-w3}).  The normal-ordering
ambiguities in $H_2$ produce the background charge
$Q \rho \cdot \partial^2\phi$ in equation~\eqref{eq:miura-T-explicit},
and the quantum correction to $H_3$ produces the composite field
$\Lambda = \normord{TT} - \frac{3}{10}\partial^2 T$.

\emph{Bar complex interpretation.}
The Toda integrability condition $\{H_2, H_3\} = 0$ is equivalent
to the associativity of the $\mathcal{W}_3$ OPE algebra, which in turn
is encoded by the nilpotence $d^2 = 0$ of the bar differential at
$c = 0$ (the uncurved locus).  The Koszul duality
$\mathcal{W}_3^k \leftrightarrow \mathcal{W}_3^{-k-6}$ exchanges
$H_2(k) \leftrightarrow H_2(-k-6)$ and $H_3(k) \leftrightarrow -H_3(-k-6)$
(the sign flip on $H_3$ reflects the orientation reversal on the Toda
phase space under $Q \mapsto -Q$).
\end{computation}

\subsection{$\mathcal{W}$-algebra bar complex dimension tables}

\begin{computation}[Bar complex dimensions across W-algebra types]
\label{comp:w-bar-dims}
\index{bar complex!dimension table!W-algebras}

We collect the bar complex dimensions for the principal W-algebras
$\mathcal{W}^k(\mathfrak{g}, f_{\mathrm{prin}})$ through degree 3.
The generators have spins $d_1+1, \ldots, d_r+1$ where $d_i$ are the
exponents of $\mathfrak{g}$.

\begin{center}
\renewcommand{\arraystretch}{1.3}
\begin{tabular}{l|ccc|cccc}
$\mathfrak{g}$ & $r$ & Exponents & Generator spins &
$\dim\bar{B}^0$ & $\dim\bar{B}^1$ & $\dim\bar{B}^2$ &
$\dim\bar{B}^3$ \\
\hline
$A_1$ & $1$ & $1$ & $2$ & $1$ & $\infty$ & $\infty$ & $\infty$ \\
$A_2$ & $2$ & $1,2$ & $2,3$ & $1$ & $\infty$ & $\infty$ & $\infty$ \\
$B_2$ & $2$ & $1,3$ & $2,4$ & $1$ & $\infty$ & $\infty$ & $\infty$ \\
$G_2$ & $2$ & $1,5$ & $2,6$ & $1$ & $\infty$ & $\infty$ & $\infty$
\end{tabular}
\end{center}
The entries are $\infty$ because the bar complex includes all modes
$L_{-n}, W_{-m}, \ldots$ with $n, m \geq 2$; the dimensions are infinite
but graded-finite by conformal weight.

\emph{Graded dimensions by conformal weight.}
For the Virasoro ($A_1$), the graded dimension of $\bar{B}^1_h$
(conformal weight $h$ states in bar degree 1) is $p_{\geq 2}(h)$
(partitions of $h$ into parts $\geq 2$; cf.\
Computation~\ref{comp:virasoro-vacuum}).
For $\mathcal{W}_3$ ($A_2$):
\begin{center}
\renewcommand{\arraystretch}{1.2}
\begin{tabular}{c|ccccccccccc}
$h$ & 2 & 3 & 4 & 5 & 6 & 7 & 8 & 9 & 10 & 11 & 12 \\
\hline
$\dim \bar{B}^1_h(\mathrm{Vir})$ &
$1$ & $1$ & $2$ & $2$ & $4$ & $4$ & $7$ & $8$ & $12$ & $14$ & $21$ \\
$\dim \bar{B}^1_h(\mathcal{W}_3)$ &
$1$ & $2$ & $3$ & $4$ & $8$ & $10$ & $17$ & $24$ & $36$ & $50$ & $76$
\end{tabular}
\end{center}
The $\mathcal{W}_3$ augmentation ideal has two strong generators
($T$ at weight 2, $W$ at weight 3), so $\bar{B}^1_h$ counts states
created by all composite fields built from $T$ and $W$ and their
derivatives at weight $h$ (modulo the vacuum).  The generating function
is:
\[
\sum_{h \geq 0} \dim \bar{B}^1_h(\mathcal{W}_3)\, q^h
= \prod_{n=2}^{\infty} \frac{1}{1-q^n}
\cdot \prod_{n=3}^{\infty} \frac{1}{1-q^n} - 1
\]
(the product of partition generating functions for the two generators,
minus the vacuum contribution).

\emph{Curvature channels.}
\begin{center}
\renewcommand{\arraystretch}{1.2}
\begin{tabular}{l|ccc|c}
$\mathfrak{g}$ & Generators & Pole orders & Curvature channels &
$c + c'$ \\
\hline
$A_1$ (Vir) & $T$ & $4$ & $1$ ($m_0 = c/2$) & $26$ \\
$A_2$ ($\mathcal{W}_3$) & $T, W$ & $4, 6$ & $2$ ($c/2$, $c/3$) & $100$ \\
$B_2$ ($\mathcal{W}(B_2)$) & $T, W^{(4)}$ & $4, 8$ & $2$ ($c/2$, $c/4$) &
$2r + 4h^\vee d$ \\
$G_2$ ($\mathcal{W}(G_2)$) & $T, W^{(6)}$ & $4, 12$ & $2$ ($c/2$, $c/6$) &
$2r + 4h^\vee d$
\end{tabular}
\end{center}
The pattern (Remark~\ref{rem:wn-vacuum-channels}): each generator
$W^{(s)}$ of spin $s$ contributes a vacuum leakage channel
$m_0^{(s)} = c/(s \cdot (2s-1)/s) = c/s$ from the pole of order
$2s$.  The total curvature is the vector $(m_0^{(s_1)}, \ldots,
m_0^{(s_r)})$, vanishing simultaneously if and only if $c = 0$.
\end{computation}

The bar-complex dimensions across the simple types (Table~\ref{comp:w-bar-dims}) display a uniform pattern: at each bar degree the chain spaces are infinite-dimensional but graded-finite under the conformal weight grading, and the growth rates are governed by a single discriminant $\Delta(x)$ that is invariant under both DS reduction and Koszul duality.  This uniformity persists in the large-$N$ limit, which we now make precise.

\subsection{\texorpdfstring{$\mathcal{W}_{1+\infty}$ and the large-$N$ limit}{W(1+infinity) and the large-N limit}}
\label{sec:w-infty}
\index{W1infinity@$\mathcal{W}_{1+\infty}$|textbf}
\index{large-$N$ limit}

The algebra $\mathcal{W}_{1+\infty}$ is the universal W-algebra
containing generators of all integer spins $\{1, 2, 3, \ldots\}$.
It arises as the algebra of area-preserving diffeomorphisms of the plane
(at the classical level) and admits a one-parameter family of large-$N$
specializations.  In the present manuscript this is not the printed
MC4 endpoint: the theorematic frontier is the filtered H-level target
for the standard principal tower together with the exact residue
identities on the finite packets~$\mathcal{I}_N$.

\begin{definition}[\texorpdfstring{$\mathcal{W}_{1+\infty}$}{W_1+infinity} as colimit]
\label{def:w-infty}
\index{W1infinity@$\mathcal{W}_{1+\infty}$!colimit}
The $\mathcal{W}_{1+\infty}$ algebra at level $(k, N)$ is the
colimit of the principal W-algebras:
\[
\mathcal{W}_{1+\infty}^{(k)}
= \varinjlim_{N} \mathcal{W}^k(\mathfrak{sl}_N, f_{\mathrm{prin}})
\]
where the transition maps
$\mathcal{W}^k(\mathfrak{sl}_N) \hookrightarrow
\mathcal{W}^k(\mathfrak{sl}_{N+1})$ are the natural inclusions
that preserve the generators $W^{(2)}, \ldots, W^{(N)}$ and add a
new generator $W^{(N+1)}$ of spin $N+1$.
\end{definition}

\begin{conjecture}[Large-N inverse-limit specialization of the
standard W-infinity MC4 package;
\ClaimStatusConjectured]
\label{conj:w-infty-bar}
\index{W1infinity@$\mathcal{W}_{1+\infty}$!bar complex}
Assume the filtered H-level package of
Definition~\ref{prop:winfty-factorization-package}: a separated
complete target whose finite quotients recover the principal stages,
together with the exact stagewise identities
\[
\mathsf{C}^{\mathrm{res}}_{s,t;u;m,n}(N)
=
\mathsf{C}^{\mathrm{DS}}_{s,t;u;m,n}(N)
\]
on the seed packets~$\mathcal{I}_N$.  Then one expects a completed
large-$N$ specialization in which the bar object of
$\mathcal{W}_{1+\infty}$ is the inverse limit of the finite-$N$ bar
complexes:
\begin{equation}\label{eq:w-infty-bar}
\barB(\mathcal{W}_{1+\infty}^{(k)})
= \varprojlim_{N} \barB(\mathcal{W}^k(\mathfrak{sl}_N))
\end{equation}
and the channel-refined characteristic profile is:
\begin{equation}\label{eq:w-infty-curvature}
\widetilde{\kappa}_N
=
\left(\frac{c_N}{2},\frac{c_N}{3},\ldots,\frac{c_N}{N}\right),
\qquad c_N = \frac{N(N^2-1)}{N+k},
\end{equation}
so the scalar summation diverges harmonically and the large-$N$ object
must retain the full channel package rather than collapse to a single
curvature scalar.  At generic 't~Hooft coupling, one further expects
the bar spectral sequence to collapse at~$E_2$, and the Koszul dual to
be
$(\mathcal{W}_{1+\infty}^{(\lambda)})^!
\simeq \mathcal{W}_{1+\infty}^{(1-\lambda)}$
\textup{(}the 't~Hooft coupling involution\textup{)}.
\emph{This is not the printed MC4 frontier.}  It is a downstream
large-$N$ specialization once the filtered target and the exact finite
identity packets are in hand.
\textup{Master conjecture parent:}
Definition~\textup{\ref{prop:winfty-factorization-package}}
\textup{(}hence MC4\textup{)}.
\end{conjecture}

\begin{remark}[Computational evidence for MC4 coefficient matching]
\label{rem:mc4-computational-evidence}
\index{MC4!computational evidence}
The finite-stage OPE coefficients and Yangian auxiliary kernels
that support the exact MC4 coefficient packets are now extracted:
\begin{enumerate}[label=\textup{(\roman*)},nosep]
\item $\mathcal{W}_\infty$ side.
  All four DS OPE coefficients are extracted as explicit rational
  functions of~$c$
  (Proposition~\ref{prop:w4-ds-ope-explicit},
  Computation~\ref{comp:w4-ds-ope-extraction}):
  \[
  c_{334}^2 = \frac{42c^2(5c+22)}{(c+24)(7c+68)(3c+46)},
  \qquad
  c_{444}^2 = \frac{112c^2(2c-1)(3c+46)}{(c+24)(7c+68)(10c+197)(5c+3)},
  \]
  \[
  \mathsf{C}_{3,4;3;0,4}^2 = \frac{9}{16}\,c_{334}^2,
  \qquad
  \mathsf{C}_{3,4;4;0,3}^2 = \frac{5}{7}\,c_{334}^2.
  \]
  Both theorematic Virasoro-target identities verified.
  Thus the first higher-spin packet is no longer DS-side extraction:
  the live stage-$4$ task is the six identities
  $\mathsf{C}^{\mathrm{res}}=\mathsf{C}^{\mathrm{DS}}$ on the packet
  $\mathcal{I}_4$ against these explicit targets.
  Channel-refined modular characteristic
  $\tilde{\kappa}_N = (c/2, c/3, \ldots, c/N)$ diverges as
  $c \cdot \log N$ (harmonic).
\item Yangian side: auxiliary kernel
  $\mathsf{K}^{\mathrm{line}}_{1,2}(N) = \Lambda^2(V)$ verified for
  $N = 2, 3, 4$, and the faithful-evaluation boundary-strip packets
  $\Delta_{a,0}(4)$ and $\Delta_{a,0}(6)$ vanish on generic tensor
  lengths $2$ and $3$ in ranks $2,3,4$.
  The remaining all-rank Yangian input is therefore the one-factor
  residue computation on the fundamental line, not the first visible
  low-stage boundary packets.
\end{enumerate}
\end{remark}

\begin{remark}[Scope: higher spin gravity and minimal model holography]
\label{rem:w-infty-holography}
\index{higher spin gravity}
\index{Gaberdiel--Gopakumar duality}
\index{minimal model holography}
The Gaberdiel--Gopakumar duality \cite{GG11} conjecturally places
$\mathcal{W}_{1+\infty}$ as the boundary algebra of Vasiliev
higher-spin gravity on $\mathrm{AdS}_3$.  In that conjectural
dictionary, the 't~Hooft coupling $\lambda$ parametrizes the bulk
gravitational coupling constant, and the Koszul duality involution
$\lambda \mapsto 1 - \lambda$ corresponds to
the exchange of the two boundary components of $\mathrm{AdS}_3$
(the bulk/boundary interchange).

The $\mathcal{W}_N$ minimal models $\mathcal{W}_N(N, N{+}1)$ at
central charge $c_N^{\min} = (N-1)(1 - N(N+1)/(N(N+1)))$ are
conjectured to be Koszul dual to $\mathcal{W}_N$ at
$c_N' = (N-1)(1 - N(N+1)/(N(N-1)))$:
\[
\mathcal{W}_N^{c_N^{\min}} \;\longleftrightarrow\;
\mathcal{W}_N^{c_N'},
\qquad c_N^{\min} + c_N' = 2(N-1) + 4N \cdot \dim(\mathfrak{sl}_N).
\]
As $N \to \infty$, this physics conjecture points toward a large-$N$
duality between the
two asymptotic regions of the higher-spin bulk, providing a
bar-cobar template for the Gaberdiel--Gopakumar conjecture.

This is a physics conjecture: a proof would require constructing
the bar complex of $\mathcal{W}_{1+\infty}$ as a completed
inverse limit, verifying the spectral sequence convergence in the
$N \to \infty$ limit, and establishing the $E_2$ collapse for
all sufficiently generic $\lambda$.  The finite-stage curved bar maps and the transition maps are themselves required inputs.
Theorem~\ref{thm:ds-koszul-intertwine} proves an associative module comparison only after the subregular $C_2$ quotient.
The full principal limit requires its own collision maps, topology, and convergence proof.

\emph{Homotopy template} (Convention~\ref{conv:hms-levels}): Type~I (existence).
The conjecture posits the existence of the bar complex of $\mathcal{W}_{1+\infty}$ as a completed inverse limit, with Koszul dual determined by the 't~Hooft coupling involution $\lambda \mapsto 1 - \lambda$.
\end{remark}


We develop the factorization Koszul dual of the $\mathcal{W}_\infty$
tower by working through the rank-$2$ case in full detail before
stating the general result.  The reader who sees why the argument works
for $\mathcal{W}_3$ will see that each ingredient generalizes to
$W_N$ without new ideas, and that the passage to $N = \infty$
adds only one further input (Mittag--Leffler for the inverse limit).

\subsubsection*{Warm-up: the Koszul dual of $\mathcal{W}_3$}

Consider the principal $\mathcal{W}$-algebra
$\mathcal{W}_3 = \mathcal{W}^k(\mathfrak{sl}_3, f_{\mathrm{prin}})$
at generic level~$k$.
It has two strong generators: the stress tensor~$T$ of conformal
weight~$2$, and the cubic Casimir field~$W$ of weight~$3$.

\emph{Step~1: the bar complex is sectorwise finite.}

A bar monomial in degree~$n$ is a tensor
$[A_1|\cdots|A_n] \otimes \omega$
where each $A_i$ is a non-vacuum state in the augmentation ideal of
$\mathcal{W}_3$ and $\omega$ is an $(n{-}1)$-form on the FM
compactification $\overline{C}_n(\bC)$.  The conformal weight
$\operatorname{wt}(A_1) + \cdots + \operatorname{wt}(A_n) = h$
defines a non-negative grading, and the key observation is:
\emph{at each weight~$h$, only finitely many monomials appear.}

The strong generators have weights in $\{2,3\}$, and a monomial of
weight~$h$ in degree~$n$ corresponds to a composition of~$h$ into~$n$
parts from $\{2,3\}$ (the weights of the strong generators and their
derivatives, with descendants at higher weights contributing only more
options at the same weight).  The dimension table
(Computation~\ref{comp:w-bar-dims}) begins:
\begin{center}
\renewcommand{\arraystretch}{1.1}
\begin{tabular}{c|ccccccccccc}
$h$ & 2 & 3 & 4 & 5 & 6 & 7 & 8 & 9 & 10 & 11 & 12 \\
\hline
$\dim \barB^1_h(\mathcal{W}_3)$ &
$1$ & $2$ & $3$ & $4$ & $8$ & $10$ & $17$ & $24$ & $36$ & $50$ & $76$
\end{tabular}
\end{center}
At weight $h = 2$: one state ($T$).
At $h = 3$: two states ($W$ and $\partial T$).
At $h = 6$, degree~$3$: the only possibility is $[T|T|T]$, with a
$2$-dimensional space of $2$-forms on $\overline{C}_3(\bC)$
(basis $\{\eta_{12} \wedge \eta_{13},\;\eta_{13} \wedge \eta_{23}\}$;
the Arnold relation kills the third).  So $\dim \barB^3_6 = 2$.

The finite-dimensionality at each weight is the input that makes the
bar complex computable: the differential $d$ is a finite matrix
in each sector, and the cohomology $H^*(\barB_h)$ is determined by
finite linear algebra.  The degree-$3$ cohomology vanishes
at every weight through $h = 16$
(Computation~\ref{comp:w3-deg3-cohom}):
\begin{center}
\renewcommand{\arraystretch}{1.1}
\begin{tabular}{c|ccccccccccc}
$h$ & 6 & 7 & 8 & 9 & 10 & 11 & 12 & 13 & 14 & 15 & 16 \\
\hline
$\dim \barB^3_h$ & 2 & 12 & 42 & 112 & 270 & 600 & 1260 & 2520 & 4884 & 9152 & 16758 \\
$\dim H^3_h$ & 0 & 0 & 0 & 0 & 0 & 0 & 0 & 0 & 0 & 0 & 0
\end{tabular}
\end{center}
The bar complex is enormous but exact in degree~$3$: this is the
Koszul property, sector by sector.  For the global statement,
$\Einf$-Koszulness of $\mathcal{W}_3$ at generic~$k$ is
Theorem~\ref{thm:master-pbw} (MC1 for principal $\mathcal{W}$-algebras).
Combined with sectorwise finiteness, this gives the factorization
bar-cobar quasi-isomorphism
$\Omega^{\mathrm{fact}}(\barB^{\mathrm{fact}}(\mathcal{W}_3))
\xrightarrow{\sim} \mathcal{W}_3$
on $\operatorname{Ran}(X)$.

\emph{Orbit transpose and the vertex carrier.}
For $\mathfrak{sl}_3$, partition transpose gives
\begin{equation}\label{eq:w3-kd-explicit}
 (3)^t=(1,1,1).
\end{equation}
Reduction at the zero nilpotent is the identity operation on the affine algebra.
This combinatorial fact does not identify the principal W-algebra with an affine vertex algebra.
Their weight-one spaces have dimensions zero and eight, respectively
(Corollary~\ref{cor:cbc28-principal-affine}).
A proposed curved Koszul identification requires a separately defined dual carrier and a map to it.

The actual subregular comparison is the inclusion into the full affine--ghost complex in Theorem~\ref{thm:cbc27-subregular-map}.
Its $C_2$ quotient gives the associative bar map of Proposition~\ref{prop:ds-koszul-hierarchy}.
On the unreduced states, the nonzero associator~\eqref{eq:cbc28-associator} rules out replacing the collision operations by ordinary normal multiplication.
Thus residue identities must follow from the actual mode maps, not from orbit transpose or equality of dimensions.

\subsubsection*{Orbit pairs in rank four}

The orbit pairs are
\[
 (4)\longleftrightarrow(1^4),\qquad
 (3,1)\longleftrightarrow(2,1,1),\qquad
 (2,2)\longleftrightarrow(2,2).
\]
Proposition~\ref{prop:cbc28-sl4-centralizers} gives the corresponding generator weights.
The principal and zero reductions again have different weight-one dimensions, now zero and fifteen.
The orbit pairing alone therefore supplies neither a conformal isomorphism nor a reflected bar map.

\subsubsection*{General theorem}

\begin{proposition}[Obstructions for a principal reduction tower;
\ClaimStatusProvedHere]\label{thm:winfty-factorization-kd}
Let $W_N^k=\mathcal W^k(\mathfrak{sl}_N,f_{\mathrm{prin}})$ be universal at noncritical level.
The orbit transpose $(N)^t=(1^N)$ does not give a conformal isomorphism from $W_N^k$ to a universal affine algebra.
Moreover, transition maps that delete the highest strong generator cannot define the full principal tower at all parameters.
\end{proposition}
\begin{proof}
The first assertion follows from the weight-one dimensions in Corollary~\ref{cor:cbc28-principal-affine}.
For the second, the $N=3$ transition would kill $W$.
At central charge $-30$, Proposition~\ref{prop:cbc28-truncation} excludes this map by its vacuum product.
\end{proof}

\begin{remark}[A completed principal comparison]
\label{eq:winfty-stage-kd}
A construction of $\varprojlim_N B^{\mathrm{ch}}(W_N)$ requires actual transition maps, their coefficient rings, and compatibility with all collision operations.
A bar--cobar comparison for that limit also requires a justified totalization and a proof controlling derived inverse limits.
The proposed reflected target $V^{-k-2N}(\mathfrak{sl}_N)$ requires a map from a defined curved Koszul dual.
Neither partition transpose nor deletion of the highest generator constructs these data.
The two proved obstructions above specify failures that an alternative construction must avoid.
\end{remark}

\begin{remark}[Reduction and bar operations]
\label{rem:ds-is-the-duality-datum}
A BRST resolution and a bar construction have different defining operations.
The former uses an odd derivation of the affine--ghost vertex algebra; the latter must specify multiplication, collision coefficients, and its coalgebra differential.
For the subregular algebra, the full vertex inclusion is Theorem~\ref{thm:cbc27-subregular-map}, while the associative quotient and its module bars are Theorems~\ref{thm:cbc28-c2-reduction} and~\ref{thm:ds-koszul-intertwine}.
The nonzero normal-product associator explains why these two constructions cannot be identified before taking the quotient.
Neither orbit transpose nor the affine level reflection supplies the missing curved coalgebra map.
\end{remark}

\subsection{DS reduction and the vacuum products}
\label{sec:ds-bar-intertwining}

\begin{computation}[A scalar test for a proposed principal comparison]
\label{comp:ds-bar-sl3-w3}
For the principal $\mathfrak{sl}_3$ reduction, the DS charge is
$c_{\mathrm{DS}}=2-24(k+2)^2/(k+3)$, whereas the affine Sugawara charge is $c_{\mathrm{aff}}=8k/(k+3)$.
At $k=0$, these are $-30$ and $0$.
Consequently a rule that sends the affine Sugawara field to the principal conformal field cannot preserve its third product:
\[
 (S_{\mathrm{aff}})_{(3)}S_{\mathrm{aff}}=0,\qquad
 (T_{\mathrm{DS}})_{(3)}T_{\mathrm{DS}}=-15\mathbf1.
\]
A unital vertex map preserving these named fields would force $0=-15\mathbf1$.
An actual DS construction uses ghosts and a conformal improvement; it cannot replace these data by a projection of affine currents onto the principal strong generators.
For subregular $\mathfrak{sl}_3$, the precise improvement and its BRST homotopy are~\eqref{eq:cbc27-stress-homotopy}.
The affine-chart collision map of Proposition~\ref{prop:cbc27-binary} then preserves every individual mode product.
This modewise statement has an explicit carrier and does not assert a completed bar comparison.
\end{computation}

\subsection{\texorpdfstring{$\mathcal{W}_3$ bar complex: degree-3 computation}{W3 bar complex: degree-3 computation}}
\label{sec:w3-bar-degree3}
\index{W3 algebra@$\mathcal{W}_3$!bar complex!degree 3}

By the sectorwise finiteness established in Step~1 above
(Computation~\ref{comp:w-bar-dims}), the degree-3 bar complex
$\barB^3_h(\mathcal{W}_3)$ is finite-dimensional at each
conformal weight~$h$.  We now compute its differential explicitly.

\begin{computation}[\texorpdfstring{$\mathcal{W}_3$}{W_3} degree-3 bar complex structure]
\label{comp:w3-deg3-structure}

At the lowest conformal weight $h = 6$ (where
$\operatorname{wt}(A) = \operatorname{wt}(B) = \operatorname{wt}(C) = 2$),
only the element $[T|T|T]$ appears.  The two form directions give:
\begin{align}
\Xi_1 &= [T|T|T] \otimes \eta_{12} \wedge \eta_{13}
\label{eq:w3-deg3-xi1} \\
\Xi_2 &= [T|T|T] \otimes \eta_{13} \wedge \eta_{23}
\label{eq:w3-deg3-xi2}
\end{align}

\emph{Differential of $\Xi_1$.}
The degree-3 bar differential has three collision terms:
\[
d(\Xi_1) = \Res_{D_{12}} + \Res_{D_{13}} + \Res_{D_{23}}
\]
Each residue extracts the OPE data at the corresponding boundary divisor.

\emph{Collision $D_{12}$ (positions 1 and 2 collide):}
\begin{align*}
\Res_{D_{12}}\Xi_1
&= \Res_{z_1 = z_2}\bigl[T(z_1) \otimes T(z_2)
\otimes T(z_3) \otimes \eta_{12} \wedge \eta_{13}\bigr] \\
&= \sum_{n \geq 0} [T_{(n)}T \,|\, T] \otimes
\Res_{z_1 = z_2}\frac{\eta_{13}}{(z_1-z_2)^{n+1}}
\end{align*}
Using the Virasoro OPE
$T(z)T(w) \sim \frac{c/2}{(z-w)^4} + \frac{2T(w)}{(z-w)^2}
+ \frac{\partial T(w)}{z-w}$:
\begin{itemize}
\item $n = 3$: $T_{(3)}T = c/2$, but $\Res_{z_1=z_2}
\frac{\eta_{13}}{(z_1-z_2)^4} = 0$ (the form $\eta_{13}
= d\log(z_1-z_3)$ has no pole at $z_1 = z_2$).
\item $n = 1$: $T_{(1)}T = 2T$, with $\Res_{z_1=z_2}
\frac{\eta_{13}}{(z_1-z_2)^2}$; expanding
$\eta_{13} = d\log(z_1-z_3) = d\log((z_2-z_3)+\epsilon)$ where
$\epsilon = z_1-z_2$, the residue extracts the $\epsilon^1$
coefficient, giving $\frac{-1}{(z_2-z_3)^2}\,dz_3 = -\wp_{23}\,dz_3$
(where $\wp_{ij} = 1/(z_i-z_j)^2$).  Thus this term contributes
$[2T|T] \otimes (-\wp_{23}\,dz_3)$.
\item $n = 0$: $T_{(0)}T = \partial T$, with $\Res_{z_1=z_2}
\frac{\eta_{13}}{z_1-z_2} = \frac{dz_2}{z_2-z_3} = \eta_{23}$.
This contributes $[\partial T|T] \otimes \eta_{23}$.
\end{itemize}

\emph{Collision $D_{13}$ (positions 1 and 3 collide):}
The form $\eta_{12} \wedge \eta_{13}$ has a pole along $D_{13}$
from the $\eta_{13}$ factor.  Setting $z_1 = z_3 + \epsilon'$:
\begin{align*}
\Res_{D_{13}}\Xi_1
&= [T|T_{(0)}T] \otimes \Res_{z_1=z_3}\eta_{12}
= [T|\partial T] \otimes \frac{dz_2}{z_2-z_3}
= [T|\partial T] \otimes \eta_{23}
\end{align*}
(only the simple pole $n = 0$ contributes, since $\eta_{12}
= d\log(z_1-z_2)$ is regular at $z_1 = z_3$).

\emph{Collision $D_{23}$:}
The form $\eta_{12} \wedge \eta_{13}$ has no pole along $D_{23}$
(neither $\eta_{12}$ nor $\eta_{13}$ involves $z_2 - z_3$ directly).
However, via the Arnold relation, we can rewrite:
\[
\eta_{12} \wedge \eta_{13}
= -\eta_{12} \wedge \eta_{23} + \eta_{13} \wedge \eta_{23}
\]
Neither summand has a pole at $D_{23}$ in the $\eta_{23}$ factor
(the $\eta_{23}$ factor is the \emph{integration variable},
not the pole source).  Hence $\Res_{D_{23}}\Xi_1 = 0$.

\emph{Result.}
\begin{equation}\label{eq:w3-dxi1}
d(\Xi_1) = [\partial T|T] \otimes \eta_{23}
+ [T|\partial T] \otimes \eta_{23}
+ [2T|T] \otimes (-\wp_{23}\,dz_3)
\end{equation}
The first two terms are degree-2 bar elements involving $\partial T$
(a descendant); the third involves a non-logarithmic 2-form
$\wp_{23}\,dz_3$ that lies in the extended form algebra
on $\overline{C}_2(\bC)$.
\end{computation}

\begin{computation}[\texorpdfstring{$\mathcal{W}_3$}{W_3} degree 3: mixed \texorpdfstring{$T$}{T}-\texorpdfstring{$W$}{W} elements]
\label{comp:w3-deg3-mixed}
\index{W3 algebra@$\mathcal{W}_3$!bar complex!mixed elements}

At conformal weight $h = 7$, the degree-3 elements involve
one $W$ generator (weight 3) and two $T$ generators (weight 2):
\begin{align*}
[T|T|W],\quad [T|W|T],\quad [W|T|T]
\end{align*}
each tensored with a basis element of $\Omega^2(\overline{C}_3)$.

\emph{The element $[W|T|T] \otimes \eta_{12} \wedge \eta_{13}$.}
The collision $D_{12}$ requires the $W(z)T(w)$ OPE:
\[
W(z)T(w) \sim \frac{3W(w)}{(z-w)^2} + \frac{\partial W(w)}{z-w}
\]
($W$ is primary of weight 3 under Virasoro).  Therefore:
\begin{align*}
\Res_{D_{12}} &= [3W|T] \otimes (-\wp_{23}\,dz_3)
+ [\partial W|T] \otimes \eta_{23}
\end{align*}

The collision $D_{13}$ requires the $T(z)T(w)$ OPE applied to
positions 1 and 3 (with $W$ at position 2 acting as spectator):
\begin{align*}
\Res_{D_{13}} &= [W|\partial T] \otimes \eta_{23}
\end{align*}

\emph{The element $[T|W|T] \otimes \eta_{12} \wedge \eta_{13}$.}
The collision $D_{12}$ uses the $T(z)W(w)$ OPE.  By
skew-symmetry of the vertex algebra OPE:
\[
T_{(n)}W = \sum_{j \geq 0} (-1)^{n+j+1}\partial^{(j)}(W_{(n+j)}T)
\]
For $n = 1$: $T_{(1)}W = 3W$ (the eigenvalue of $L_0$ on $W$).
For $n = 0$: $T_{(0)}W = \partial W$.  Therefore:
\begin{align*}
\Res_{D_{12}} &= [3W|T] \otimes (-\wp_{23}\,dz_3)
+ [\partial W|T] \otimes \eta_{23}
\end{align*}
The collision $D_{13}$ uses $T(z)T(w)$ again:
\begin{align*}
\Res_{D_{13}} &= [W|2T] \otimes (-\wp_{23}\,dz_3)
+ [W|\partial T] \otimes \eta_{23}
\end{align*}
\end{computation}

\begin{computation}[Arnold cancellation at \texorpdfstring{$\mathcal{W}_3$}{W_3} degree 3]
\label{comp:w3-arnold-deg3}
\index{Arnold relation!$\mathcal{W}_3$ degree 3}

The key structural feature at degree~3 is the cancellation of
vacuum leakage.  Consider the degree-3 element
$\Xi = [T|T|T] \otimes (\eta_{12} \wedge \eta_{23})$ and compute
$d(\Xi)$.

\emph{Vacuum channel.}
The potential vacuum contribution from $D_{12}$ is
$T_{(3)}T = c/2$ (the quartic pole).  This requires extracting the
residue of $c/2 \cdot T(z_3) \cdot \eta_{12} \wedge \eta_{23}$
at $z_1 = z_2$.  Since $\eta_{12} \wedge \eta_{23}$ has a
simple pole at $D_{12}$ from the $\eta_{12}$ factor:
\[
\Res_{D_{12}}\bigl[\tfrac{c/2}{(z_1-z_2)^4}
\cdot \eta_{12} \wedge \eta_{23}\bigr]
= \tfrac{c/2}{3!} \cdot
\partial_\epsilon^3\bigl|_{\epsilon=0}
\bigl[\tfrac{d\epsilon}{\epsilon} \wedge \eta_{23}\bigr] = 0
\]
(the form $d\epsilon/\epsilon$ contributes a simple pole, but the
$1/(z_1-z_2)^4$ factor requires a fourth-order pole; the product
has only a $1/\epsilon^3$ singularity, and
$\Res_{\epsilon=0}[\epsilon^{-3}\,d\epsilon] = 0$).

\emph{Total vacuum leakage.}
By the same mechanism at $D_{13}$ and $D_{23}$: each collision
produces a residue involving $c/2$ times a form with insufficient
pole order.  The vacuum leakage at degree~3 therefore vanishes:
\begin{equation}\label{eq:w3-deg3-no-vacuum}
\bigl(d(\Xi)\bigr)\big|_{\text{vacuum}} = 0
\end{equation}
This is the $\mathcal{W}_3$ analog of
Proposition~\ref{prop:arnold-virasoro-deg3}: the Arnold relation
kills the vacuum channel at degree~3, ensuring that the bar
differential does not map degree-3 elements to the vacuum
(which would signal a failure of the filtration).
\end{computation}

\begin{proposition}[\texorpdfstring{$\mathcal{W}_3$}{W_3} degree-3 vacuum cancellation; \ClaimStatusProvedHere]
\label{prop:w3-deg3-vacuum}
\index{vacuum leakage!$\mathcal{W}_3$ degree 3}
In the degree-$3$ bar differential of the $\mathcal{W}_3$ algebra,
no element maps to the vacuum $|0\rangle \in \bar{B}^0$.
Equivalently, the composition $\bar{B}^3 \xrightarrow{d}
\bar{B}^2 \xrightarrow{d} \bar{B}^1 \xrightarrow{d}
\bar{B}^0$ maps $\bar{B}^3$ to zero in $\bar{B}^0$.
\end{proposition}

\begin{proof}
The only way a degree-3 element could map to the vacuum is through
iterated residues: $d \circ d|_{\bar{B}^3 \to \bar{B}^0}$.
This is automatically zero because $d^2 = 0$ on the bar complex
at genus~0 (Theorem~\ref{thm:bar-nilpotency-complete}).  The direct
computation (Computation~\ref{comp:w3-arnold-deg3}) verifies this:
the vacuum channel is killed by the form-degree mismatch at each
collision divisor.
\end{proof}


\subsection{\texorpdfstring{$\mathcal{W}_3$ degree-3 cohomology}{W3 degree-3 cohomology}}

\begin{computation}[\texorpdfstring{$\mathcal{W}_3$}{W_3} degree-3 cohomology at low weight]
\label{comp:w3-deg3-cohom}
\index{W3 algebra@$\mathcal{W}_3$!bar cohomology!degree 3}

At conformal weight $h = 6$ (the lowest weight for degree-3
elements), the bar complex is:
\[
\bar{B}^3_6 = \bC\{[T|T|T]\} \otimes \Omega^2(\overline{C}_3)
\cong \bC^2
\]
The differential $d \colon \bar{B}^3_6 \to \bar{B}^2_6$ is
computed in Computation~\ref{comp:w3-deg3-structure}.  Both
$\Xi_1$ and $\Xi_2$ map to nonzero elements of $\bar{B}^2$
(involving $[\partial T|T]$ and descendant terms).

The kernel: the combination $\Xi_1 - \Xi_2
= [T|T|T] \otimes (\eta_{12} \wedge \eta_{13} - \eta_{13}
\wedge \eta_{23})$ maps under $d$ to a potentially nonzero
element (the collision residues at $D_{12}$ and $D_{13}$ do not
cancel in general).  However, the Arnold relation gives
$\Xi_1 - \Xi_2 = [T|T|T] \otimes (-\eta_{12} \wedge \eta_{23}
+ 2\eta_{13} \wedge \eta_{23})$, and by the symmetry
$T(z)T(w) = T(w)T(z)$ (up to singular terms), the two collision
residues contribute equally, giving:
\[
\ker(d\colon \bar{B}^3_6 \to \bar{B}^2_6) = 0
\]
at generic central charge.  Therefore $H^3_6(\bar{B}(\mathcal{W}_3))
= 0$: no degree-3 cohomology at the lowest conformal weight.

At weight $h = 7$, the degree-3 space has dimension~$12$: three
orderings of strong generators $(T, T, W)$, each with $\dim \bar{V}_3 = 2$
states at weight~$3$ (namely $W$ and $\partial T$), times
two form directions.  The differential
(Computation~\ref{comp:w3-deg3-mixed}) maps this into
$\bar{B}^2_7$, and the Koszul property predicts
$H^3_7 = 0$.

\emph{Dimension table through weight $h = 16$.}
\begin{center}
\renewcommand{\arraystretch}{1.2}
\begin{tabular}{c|ccccccccccc}
$h$ & 6 & 7 & 8 & 9 & 10 & 11 & 12 & 13 & 14 & 15 & 16 \\
\hline
$\dim \bar{B}^3_h$ & 2 & 12 & 42 & 112 & 270 & 600 & 1260 & 2520 & 4884 & 9152 & 16758 \\
$\dim H^3_h$ & 0 & 0 & 0 & 0 & 0 & 0 & 0 & 0 & 0 & 0 & 0
\end{tabular}
\end{center}
The vanishing $H^3_h = 0$ for $h \leq 16$ is consistent with the
Koszul property of the $\mathcal{W}_3$ algebra
(Theorem~\ref{thm:w-algebra-koszul}).
\end{computation}


\subsection{\texorpdfstring{$\mathcal{W}_3$ bar complex at the dual level}{W3 bar complex at the dual level}}
\label{sec:w3-dual-level-bar}
\index{W3 algebra@$\mathcal{W}_3$!dual level}

\begin{computation}[\texorpdfstring{$\mathcal{W}_3$}{W_3} curvature at the dual level]
\label{comp:w3-curvature-dual}
\index{W3 algebra@$\mathcal{W}_3$!curvature!dual level}

Under the Feigin--Frenkel involution $k \mapsto k' = -k-6$
(Corollary~\ref{cor:level-shifting-part1}):
\begin{center}
\renewcommand{\arraystretch}{1.3}
\begin{tabular}{l|cc}
& Level $k$ & Dual level $k' = -k-6$ \\
\hline
Central charge & $c = 2 - \frac{24(k+2)^2}{k+3}$ &
$c' = 2 - \frac{24(k'+2)^2}{k'+3} = 100 - c$ \\[6pt]
$T$-curvature & $m_0^{(T)} = c/2$ &
$(m_0^{(T)})' = c'/2 = (100-c)/2$ \\[4pt]
$W$-curvature & $m_0^{(W)} = c/3$ &
$(m_0^{(W)})' = c'/3 = (100-c)/3$ \\[4pt]
Sum & $m_0^{(T)} + (m_0^{(T)})' = 50$ &
$m_0^{(W)} + (m_0^{(W)})' = 100/3$
\end{tabular}
\end{center}

\emph{Special levels.}
\begin{itemize}
\item $k = -3$ (critical level): $c$ has a pole ($c \to \pm\infty$
depending on direction of approach).  The Sugawara construction
degenerates, and the $\mathcal{W}_3$ algebra at critical level is
the center of the completed universal enveloping algebra
$\widehat{U}(\widehat{\mathfrak{sl}}_{3,-3})$.
\item $k = 0$: $c = 2 - 24 \cdot 4/3 = -30$,
$c' = 130$, $m_0^{(T)} = -15$, $m_0^{(W)} = -10$.
\item $k = 1$: $c = 2 - 24 \cdot 9/4 = -52$,
$c' = 152$, $m_0^{(T)} = -26$, $m_0^{(W)} = -52/3$.
\item $k = -5$ ($k' = -1$): $c = 2 - 24 \cdot 9/(-2) = 110$,
$c' = -10$, $m_0^{(T)} = 55$, $m_0^{(W)} = 110/3$.
\end{itemize}

\emph{Zero locus.}
Both curvature channels vanish simultaneously if and only if $c = 0$:
\[
c = 0 \iff 12k^2 + 47k + 45 = 0 \iff k = -\tfrac{9}{4},\; -\tfrac{5}{3}
\quad (\text{non-integral rational levels: } k = -\tfrac{9}{4},\; -\tfrac{5}{3})
\]
The zero-curvature locus consists of two rational non-integral levels, in contrast
to the affine Kac--Moody case where $m_0 = 0$ at the critical level
$k = -h^\vee$ (a single rational value).  At these levels, the
$\mathcal{W}_3$ bar complex is strict ($d^2 = 0$ at all genera)
and the Koszul duality is of classical type.
In the language of Theorem~\ref{thm:explicit-theta}, the zero locus
is precisely where the universal Maurer--Cartan class vanishes:
$\Theta_{\mathcal{W}_3^k} = 0$ iff $c = 0$ iff
$k \in \{-9/4, -5/3\}$.  Away from these values,
$\Theta_{\mathcal{W}_3^k} = \kappa(\mathcal{W}_3^k)
\cdot \eta \otimes \Lambda$ with
$\kappa = c/2 = 1 - 12(k+2)^2/(k+3)$
(for the $T$-channel) controls the curvature at every genus.
\end{computation}

\begin{computation}[Subregular binary vacuum products]
\label{comp:bp-bar}
\index{Bershadsky--Polyakov algebra!binary collisions}
For $k\ne-3$, the exact products in Proposition~\ref{prop:cbc27-bp-products} give
\begin{align}
 J_{(1)}J&=\frac{2k+3}{3}\mathbf1,\label{eq:bp-curv-sl2}\\
 G^+{}_{(2)}G^-&=(k+1)(2k+3)\mathbf1,\label{eq:bp-curv-G}\\
 T_{(3)}T&=-\frac{(2k+3)(3k+1)}{2(k+3)}\mathbf1.
\end{align}
The normalization $[a_\lambda b]=\sum_{n\geq0}\lambda^na_{(n)}b/n!$
accounts for the factor $2$ between the $\lambda^2$ coefficient and the cubic pole.
All three scalar terms vanish at $k=-3/2$.
At $k=-1/2$ their values are $2/3,1,1/5$, respectively.
Thus that level still has nonzero vacuum collisions.

At $k=-3/2$, the universal algebra still has the four independent PBW generators.
For example, $J_{(0)}G^+=G^+$ remains nonzero.
Vanishing scalar terms therefore do not imply a zero collision operation.
The polynomial representatives $J,U,P,V$ and their strict reduction map remain defined at $k=-3$.
The conformal formula for $T$ requires inversion of $k+3$ and has a different domain.
\end{computation}


\subsection{$\mathcal{W}$-algebra Koszul duality across the DS hierarchy}
\label{sec:ds-hierarchy-koszul}
\index{Drinfeld--Sokolov reduction!Koszul duality hierarchy}

\begin{proposition}[Subregular reduction and associative bars;
\ClaimStatusProvedHere]\label{prop:ds-koszul-hierarchy}
For the universal subregular $\mathfrak{sl}_3$ family over $\mathbb C[k]$, the reduction induces
\[
 B(R_{W_R})\xrightarrow{\ B(i)\ }B(R_{C_R})
\]
as a quasi-isomorphism of direct-sum differential coalgebras with the augmentations of Theorem~\ref{thm:cbc28-bar-comparison}.
Their cohomology is the exterior algebra on four classes of degree $-1$.
On the full vertex carrier, the normal-product bar already fails at length three:
\[
 b_2^2(sJ\otimes sJ\otimes sU)=2s\bigl(:(\partial J)U:\bigr)\ne0.
\]
\end{proposition}
\begin{proof}
The polynomial inclusion, projection, and contraction are those of Theorem~\ref{thm:cbc28-c2-reduction}.
Theorem~\ref{thm:cbc28-bar-comparison} extends the two maps letterwise to the associative bars.
Corollary~\ref{cor:cbc28-bar-cohomology} computes their cohomology.
The final identity is Proposition~\ref{prop:cbc28-bar-obstruction} and holds before taking the quotient.
\end{proof}

\begin{remark}[Orbit transpose and a reflected chiral target]
For a nilpotent $f$ of a simple Lie algebra, the proposed relation
\[
 (\mathcal W^k(\mathfrak g,f))^!\simeq
 \mathcal W^{-k-2h^\vee}(\mathfrak g^\vee,f^{\mathrm{BV}})
\]
requires a specified curved chiral dual and a map between the two coefficient systems.
The same-level reduction and its $C_2$ quotient do not define this cross-level map.
In the subregular $\mathfrak{sl}_3$ case, the vacuum coefficients prevent augmentations on both reflected factors, and the normal product fails associativity.
These are explicit equations that any proposed extension must address.
\end{remark}

\begin{computation}[Nilpotent reductions for $\mathfrak{sl}_4$]
\label{comp:sl4-ds-hierarchy}
\index{sl4@$\mathfrak{sl}_4$!nilpotent orbits}
Proposition~\ref{prop:cbc28-sl4-centralizers} computes the orbit dimensions and strong generator weights for all five partitions.
For the subregular partition $(3,1)$, the centralizer has dimension five and generator weights $1,2,2,2,3$.
For the minimal partition $(2,1,1)$, the centralizer has dimension nine and weights $1^4,(3/2)^4,2$.
All fields are even.

Partition transpose gives the orbit pairs
\[
 (1^4)\longleftrightarrow(4),\qquad
 (2,1,1)\longleftrightarrow(3,1),\qquad
 (2,2)\longleftrightarrow(2,2).
\]
The corresponding universal vertex algebras have different generator weights.
A vertex comparison requires a map and its preserved structures in addition to this orbit pairing.
For the principal and zero orbits, Corollary~\ref{cor:cbc28-principal-affine} excludes a conformal isomorphism.

If the generator weights, counted with multiplicity, are $d_1,\ldots,d_m$, their universal vacuum series is
\[
 \prod_{i=1}^m\prod_{r\geq0}(1-q^{d_i+r})^{-1}.
\]
This follows from the ordered PBW monomials.
It counts states rather than tensor-bar chain groups or their cohomology.
\end{computation}
